\documentclass[11pt,a4paper]{article}
\usepackage[utf8]{inputenc}
\usepackage[bulgarian]{babel}
\usepackage{amsmath}
\usepackage{amsfonts}
\usepackage{amssymb}
\usepackage{graphicx}
\usepackage{geometry}
\geometry{margin=2.5cm}

\title{Архитектурна Спецификация на MORT: Микро-оптотронен Реактивен Теселатор \\ \large Версия 2.0: Хибридна Вълнова Инфраструктура и Самолекуваща се Топология}
\author{Инженерен Доклад по Проект MORT}
\date{\today}

\begin{document}

\maketitle

\begin{abstract}
Настоящият доклад описва разширената архитектура на MORT (Micro-Optotronic Reactive Tesselator) – компютърна система от ново поколение, базирана на не-фон-Нойманова архитектура. Системата използва вълнова интерференция и топологична теселация за извършване на масивно паралелни изчисления със скоростта на светлината. Документът детайлизира решенията на критичните проблеми с кохерентността, синхронизацията и термичната стабилност чрез въвеждането на М.А.М. часовници, RF Backflip детектори и динамични електромагнитни щитове.
\end{abstract}

\section{Въведение в Теселационната Логика}
За разлика от традиционната бинарна логика, MORT оперира чрез „Прожектирана регенерация“. Входящите сигнали, модулирани като терахерцови вълни ($f \approx 300$ GHz), се сблъскват с топологично ядро (Shield), което деформира вълновия фронт. Резултатът от изчислението е физическата картина на интерференция (теселация), която се заснема от високоскоростна фотонна матрица.

\section{Хардуерни Модули и Решения}

\subsection{D-куплунг и Регенерация на Сигнала}
За справяне с ентропията и затихването, системата внедрява \textbf{D-куплунг} (Detector-Coupler). 
\begin{itemize}
    \item \textbf{Принцип:} Активен фотосензор-регенератор, който улавя отслабената вълна и я препредава с първоначалната ѝ амплитуда, запазвайки фазовата информация.
    \item \textbf{Функция:} Служи като „репитер“ в ПАМ (Периметърна Аритметична Памет) и като интерфейс между отделните ядра.
\end{itemize}

\subsection{М.А.М. (Monophotonic Atomic Marker)}
Проблемът със синхронизацията на пикосекундни интервали се решава чрез \textbf{М.А.М. модул}. 
\begin{itemize}
    \item \textbf{Механизъм:} Един-единствен фотон, затворен в свръхстабилен фиброоптичен пръстен, служи за квантов метроном.
    \item \textbf{Резултат:} Осигурява „сърцебиенето“ на системата, по което се синхронизират ПАМ и изчислителните ядра, елиминирайки дрейфа.
\end{itemize}

\subsection{RF Backflip Detector (Имунна система)}
За премахване на паразитното „ехо“ и обратните отражения е внедрен многочестотен \textbf{Backflip Detector}.
\begin{itemize}
    \item \textbf{Реактивност:} При засичане на нежелана обратна вълна, детекторът генерира анти-вълна (сигнален офсет), която деструктивно интерферира с шума, стабилизирайки системата в реално време.
\end{itemize}

\section{Мрежова Архитектура и ТрансБука}
За комуникация между различни MORT единици се използва хибриден модел на „Универсален транслатор“:
\begin{enumerate}
    \item \textbf{Локален Диалект:} Всяко ядро работи със свой уникален офсет, определен от специфичната му електромагнитна топология.
    \item \textbf{ТрансБука (TransBook):} Входно-изходен транслатор, който конвертира локалния „диалект“ в „Златен стандарт“ за пренос по мрежата.
    \item \textbf{Пилотна вълна:} Паралелен сигнал, който следи за изкривявания по трасето и позволява на приемника да адаптира офсетите автоматично.
\end{enumerate}

\section{Сигурност и Криптиране}
Сигурността е интегрирана физически чрез \textbf{SHA-256000000}. Алгоритъмът се изпълнява нативно чрез вълнов трансдюсер, който превръща бинарните данни в сложни аудио-визуални вълнови потоци (PNG към WAV транслация). Това прави хакването невъзможно без физическото притежание на идентична електромагнитна структура.

\section{Процес на Поддръжка: Мексиканска вълна}
Вместо пълен рестарт, MORT използва \textbf{Мексиканска вълна} за системно продухване.
\begin{itemize}
    \item \textbf{Динамичен Рефреш:} Последователен вълнов фронт преминава през всички модули, изчиствайки натрупаните грешки и калибрирайки фазата, без да прекъсва основния изчислителен процес.
\end{itemize}

\section{Енергийна ефективност и Генериране}
Чрез използването на Wolfram-базирани самоизчисляващи се алгоритми, MORT постига състояние на „енергийна затвореност“. Вълновата енергия, използвана за изчисление, се рециклира обратно в ПАМ модула. Системата функционира не само като изчислител, но и като генератор на кохерентна енергия чрез подреждане на информационната ентропия.

\section{Заключение}
MORT представлява преход от статичната електроника към динамичната полева физика. Със своята самолекуваща се структура и вълнова природа, системата предлага практически неограничена изчислителна мощ при минимален енергиен отпечатък, поставяйки началото на ерата на Когнитивния хардуер.

\end{document}
